\section{Taller resuelto: Una trayectoria integral con viento de cola}

\noindent\textbf{Nivel:} Intermedio.

\apartado{1. Enunciado}{Calcule alcance y duración del mismo planeador desde 3100 hasta 350 m ISA, manteniendo CL de mejor planeo en calma, con viento de cola 9 m/s. Compare avance terrestre medio con TAS inicial.}

\apartado{2. Datos}{S=16.2 m², masa 455 kg, polar 0.014+0.032CL²; viento vertical nulo.}

\apartado{3. Incógnitas}{Las magnitudes y comprobaciones solicitadas en el caso.}

\apartado{4. Sistema de referencia}{Ruta aérea y terrestre alineadas; altitud geopotencial.}

\apartado{5. Hipótesis físicas}{Masa y polar constantes, sucesión cuasiestacionaria de equilibrios locales.}

\apartado{6. Modelo matemático}{Una trayectoria integral con viento de cola}

\apartado{7. Ecuaciones simbólicas}{Se parte de los balances y definiciones siguientes.}

\begin{equation}R_a=\Delta hC_L/C_D\end{equation}

\begin{equation}E=\int\dd h/s(h)\end{equation}

\begin{equation}R_g=R_a+9E,\quad\bar V_{g,H}=R_g/E\end{equation}

\apartado{8. Desarrollo algebraico}{Las operaciones se realizan en este orden, conservando las condiciones del modelo.}

\begin{enumerate}

\item Se evalúa el CL de máxima eficiencia.

\item Se integra la tasa a densidad variable y se añade el desplazamiento de viento.

\item La velocidad media terrestre se obtiene como distancia dividida por tiempo y no coincide con una TAS puntual.

\end{enumerate}

\apartado{9. Sustitución numérica}{Se conservan los resultados intermedios sin redondear y se muestran todas las operaciones escalares disponibles.}

\textbf{Integral de tiempo desarrollada.} Calcular $C_D=C_{D0}+kC_L^2$ y $\delta=\arctan(C_D/C_L)$. Del equilibrio exacto, $V(h)=\sqrt{2mg_0\cos\delta/[\rho(h)SC_L]}$ y $s(h)=V(h)\sin\delta$. Definir $K=\sqrt{2mg_0\cos\delta/(SC_L)}\sin\delta$ da $s=K/\sqrt\rho$. En la troposfera, $\rho=\rho_0(1-\beta h/T_0)^a$, con $a=g_0/(R\beta)-1$.\[\fitmath{E=\frac{\sqrt{\rho_0}}{K}\int_{h_f}^{h_i}(1-\beta h/T_0)^{a/2}dh=\frac{\sqrt{\rho_0}T_0}{K\beta(1+a/2)}\left[(1-\beta h_f/T_0)^{1+a/2}-(1-\beta h_i/T_0)^{1+a/2}\right]}\]Los parámetros intermedios son $C_L=0.6614378$, $C_D=0.028$, $\delta=0.04230676\,\mathrm{rad}$ y $a=4.25588$. La evaluación por alturas permite contrastar el integrando:\begin{center}\begin{tabular}{rrrrr}
\toprule $h$ (m) & $
ho$ & $V$ (m/s) & $s$ (m/s) & $1/s$ (s/m)\\\midrule
$350$ & $1.184365$ & $26.50586$ & $1.121043$ & $0.8920268$\\
$1037.5$ & $1.107554$ & $27.40957$ & $1.159264$ & $0.8626161$\\
$1725$ & $1.034613$ & $28.3593$ & $1.199432$ & $0.8337277$\\
$2412.5$ & $0.9654061$ & $29.35821$ & $1.24168$ & $0.8053602$\\
$3100$ & $0.8997973$ & $30.40971$ & $1.286153$ & $0.7775127$\\
\bottomrule\end{tabular}\end{center}\[\fitmath{E\approx2293.706\,\mathrm s,\qquad R_a=(h_i-h_f)\frac{C_L}{C_D}\approx64962.64\,\mathrm m}\]\[\fitmath{R_g=R_a+U E\approx85606\,\mathrm m}\]





\apartado{10. Resultados con unidades}{Los resultados pertenecen al modelo y al intervalo de datos declarados.}

\begin{center}\begin{tabular}{lr}\toprule Magnitud&Resultado\\\midrule

Alcance aéreo & $64962.64\,\mathrm m$\\

Duración & $2293.706\,\mathrm s$\\

Alcance terrestre & $85606\,\mathrm m$\\

Avance terrestre medio & $37.32213\,\mathrm{m\,s^{-1}}$\\

\bottomrule\end{tabular}\end{center}

\apartado{11. Interpretación}{El viento de cola aumenta alcance sin cambiar la duración de la ley de descenso definida.}

\apartado{12. Comprobación independiente y alcance}{Comparar cuadratura con primitiva ISA y comprobar Rg-Ra=9E.}

\apartado{13. Limitaciones}{No se ha optimizado la velocidad para aprovechar el viento; sólo se evalúa la ley impuesta.}

\par\smallskip\noindent{\footnotesize\textit{Datos y solución originales. Las sustituciones registradas se reproducen con SymPy y Studium; eso no establece validación experimental ni revisión experta.}}
